\documentclass[preprint,12pt]{elsarticle}

\usepackage{amsmath,amssymb}
\usepackage{graphicx}
\graphicspath{{../figures/}}
\usepackage{booktabs}
\usepackage[section]{placeins}
\renewcommand{\topfraction}{0.9}
\renewcommand{\bottomfraction}{0.8}
\renewcommand{\textfraction}{0.05}
\renewcommand{\floatpagefraction}{0.75}
\setcounter{topnumber}{3}
\setcounter{totalnumber}{5}
\usepackage[colorlinks=true,linkcolor=blue,citecolor=blue,urlcolor=blue]{hyperref}

\journal{Journal of Quantitative Spectroscopy and Radiative Transfer}

\begin{document}

\begin{frontmatter}

\title{When does dawn become visible? The Quranic dawn (fajr al-sadiq)
and dusk (shafaq) as a detection problem, tested against 116 years of
observation}

\author[ku]{Mostafa Mahdi}
\affiliation[ku]{organization={University of Copenhagen},
            addressline={rhk510@alumni.ku.dk},
            city={Copenhagen},
            country={Denmark}}

\begin{abstract}
Islamic prayer timekeeping requires the observation, or prediction, of
four twilight events: the true dawn (fajr al-sadiq, the white thread of
light spreading along the horizon), the false dawn that precedes it (fajr
al-kadhib, the zodiacal wolf's tail), and the evening disappearance of the
red and of the white twilight (shafaq al-ahmar and al-abyad). Modern
practice replaces these observations with fixed solar depression angles,
and the major conventions disagree by tens of minutes because a fixed
angle cannot represent what the sources describe: a visual detection
event, dependent on atmosphere, background sky, and the human observer.
We formulate the dawn criterion of Quran 2:187 as a quantitative
psychophysical detection problem: simulated absolute sky radiance across
a horizon band of patches, mesopic photometry, and a Weber contrast test
of each patch's growth above its own deep-night baseline, requiring
simultaneous distinctness across the band's lateral extent. Sky radiance
is supplied by a validated backward Monte Carlo radiative transfer model
(companion paper); the detection layer contains a single calibrated
constant bridging laboratory contrast thresholds to field detection of a
borderless gradient. Calibrated on the desert campaigns that report a
documented central-mean dawn depression, the model reproduces sixteen
published visual and instrumental
campaigns spanning 116 years, latitudes 6$^\circ$S to 52$^\circ$N, and
desert, agricultural, coastal, and urban backgrounds, with a median
absolute residual of 0.26$^\circ$ of solar depression over the fourteen
scored rows. An out-of-sample
year-long panel study at Birmingham, UK, is reproduced with zero
retuning (mean residual $+0.54^{\circ}$, RMS $0.95^{\circ}$ over all
42 published dates), including its double-peaked seasonal curve and
its summer relaxation to about 12.5$^\circ$; on the same benchmark the
fixed angles in worldwide use miss by 1.8 to 4.4$^{\circ}$ RMS, and
the most widely used one is undefined on 20 of the 42 dates. The
calibrated constant inverts consistently ($f = 51$ to 57) across every
campaign that documents a central mean, with no latitude trend: it
behaves as a single property of human dawn detection rather than a fit
artifact. We further
compute the false dawn as a distinct earlier detection of the zodiacal
column, reproduce the qualitative suppression of its distinctness by
urban skyglow, and quantify how measured cloud fields, aerosol, and
artificial light shift the times, with propagated uncertainties reported
per event. Computed times are offered as a research instrument, not a
replacement for established calendars.
\end{abstract}

\begin{keyword}
fajr \sep twilight \sep prayer times \sep visual detection \sep contrast
threshold \sep zodiacal light \sep sky brightness
\end{keyword}

\end{frontmatter}

\section{Introduction}
\label{sec:intro}

The Islamic day is structured by observable celestial events, and two of
them are twilight phenomena of considerable optical subtlety. The Quran
defines the onset of the daily fast, which coincides with the earliest
time of the dawn prayer, by a visual criterion: ``eat and drink until the
white thread of dawn becomes distinct to you from the black thread [of
night]'' (Quran 2:187). Prophetic tradition sharpens the geometry: the
true dawn (al-fajr al-sadiq) ``spreads laterally along the horizon,''
while the false dawn that precedes it (al-fajr al-kadhib) rises
``vertically, like the tail of the wolf,'' and does not begin the fast.
The evening prayer window is likewise bounded by the disappearance of
twilight (al-shafaq), with the classical schools differing over whether
the red glow or the later white glow is decisive.

Contemporary practice almost universally replaces these observations with
fixed solar depression angles: $18^\circ$ (Muslim World League),
$15^\circ$ (ISNA), $19.5^\circ$ (Egyptian General Authority), or fixed
clock offsets (Umm al-Qura). The conventions disagree with one another by
tens of minutes at mid-latitudes and diverge unboundedly toward the
poles, where some depression angles are never reached in summer. The
disagreement is structural rather than empirical: a fixed angle asserts
that dawn visibility is a function of solar geometry alone, whereas the
defining texts describe a detection event, which necessarily depends on
the radiance of the dawn band, the darkness of the adjacent sky, the
state of the atmosphere, and the human visual system. Observation
campaigns confirm the structural point: measured naked-eye dawn
depressions cluster near $14$ to $15^\circ$ under desert skies
\citep{almostafa2005,khalifa2018,mawad2024}, sit lower under urban or
humid skies, and drift with season at high latitude
\citep{openfajr2016}, none of which a constant can express.

This paper takes the textual criterion seriously as a physics problem.
We formulate ``the white thread becomes distinct to you'' as a
psychophysical detection test on simulated absolute sky radiance, with
the simulation supplied by a validated backward Monte Carlo radiative
transfer model of the twilight sky (described and validated in a
companion methods paper, hereafter Paper~I). The criterion is relative
by construction, comparing tonight's dawn band against tonight's night
sky, which gives it two properties no absolute convention has: systematic
radiometric errors largely cancel, and the criterion remains defined at
high latitude, where the sky brightens without ever becoming fully dark.

Section~\ref{sec:sources} reviews the textual and observational
constraints the criterion must satisfy. Section~\ref{sec:criterion}
formulates the detection model. Section~\ref{sec:calibration} describes
its single calibrated constant and the inversion study that tests the
constant's universality. Section~\ref{sec:campaigns} presents the
validation against sixteen campaigns, including the out-of-sample
Birmingham panel year. Section~\ref{sec:kadhib} treats the false dawn.
Section~\ref{sec:inputs} quantifies the environmental response: clouds,
aerosol, artificial skyglow, and the propagated uncertainty of each
computed time. Section~\ref{sec:discussion} discusses jurisprudential
scope and limitations.

\section{Textual and observational constraints}
\label{sec:sources}

The criterion is constrained on three sides. Textually: the event is a
distinctness threshold (``tabayyana,'' becomes distinct) between the lit
band and the adjacent night sky, judged by an unaided human observer
(``lakum,'' to you); the true dawn is extended laterally
(``mustatir(an)'') while the false dawn is a narrow vertical column
(``mustatil(an),'' the wolf's tail); and the evening window closes with
the disappearance of shafaq, red or white by school.

Observationally: modern campaigns supply the geometry and the depression
statistics. The twilight arch at the moment of first distinctness sits
at approximately $2.66 \pm 0.23^\circ$ of apparent altitude
\citep{mawad2024}; the true dawn band spans roughly 30 to $40^\circ$ of
azimuth at detection \citep{ilyas1984,khalifa2018}; the false dawn's
zodiacal wedge is about $20^\circ$ wide at its base \citep{sultan2004}.
Depression statistics under dark desert skies: $14.6 \pm 0.3^\circ$
(naked eye, one year, twice monthly, eight observers;
\citealp{almostafa2005}); $14.01 \pm 0.32^\circ$ over $\sim$80 mornings
\citep{khalifa2018}; $14.90 \pm 0.17^\circ$ by calibrated camera
\citep{mawad2024}; and, remarkably, a photographic value of $14.25^\circ$
for the first color difference from Aswan in winter 1908
\citep{miethe1909}.

Psychophysically: laboratory contrast thresholds
\citep{blackwell1946,crumey2014} predict that the dawn band's photometric
excess over the night sky crosses detection threshold near depression
$17$ to $18^\circ$, several degrees before any field campaign reports
naked-eye detection. The discrepancy is not a contradiction; it is the
known gap between detecting a sharp-edged laboratory stimulus and
noticing a borderless luminance gradient spread over tens of degrees of
sky, compounded by field conditions (imperfect dark adaptation, unknown
fixation, criterion conservatism). Photometers and long-exposure cameras
do detect the earlier excess, which is precisely why instrumental
campaigns report systematically earlier (deeper) values than naked-eye
panels observing the same sky. Any credible model of the textual
criterion must bridge this gap explicitly rather than absorb it into
retuned physics.

\section{The detection criterion}
\label{sec:criterion}

For each scanned solar depression the model simulates seven sky patches
at an apparent altitude of $3^\circ$ (the scan-grid altitude nearest
the measured arch altitude of
Section~\ref{sec:sources}): five ``band'' patches spanning
$\pm 18^\circ$ of azimuth about the solar azimuth (the mustatir extent),
and two ``reference'' patches at $\pm 100^\circ$, representing the night
sky against which distinctness is judged. Each patch's radiance is the
sum of the Monte Carlo twilight field (Paper~I), the direction-dependent
celestial background (zodiacal light, integrated starlight, airglow,
moonlight), and artificial skyglow, converted to mesopic luminance from
its simulated spectrum.

Detection is a Weber contrast test with the structure of the ayah (the
verse). Let
$L_i(\theta)$ be the mesopic luminance of band patch $i$ at depression
$\theta$, and let $B_i$ be that same patch's deep-night baseline (its
own value in full night, so that the celestial background subtracts
patch by patch). The patch's signal is its growth above its own
baseline, $\Delta L_i(\theta) = L_i(\theta) - B_i$, and it is judged
against the observer's adaptation, set by the reference-sky luminance
$L_{\mathrm{ref}}$, through the threshold-versus-intensity function:
patch $i$ is distinct when
\begin{equation}
\Delta L_i(\theta) \;\ge\; f \cdot C_{\mathrm{thr}}(L_{\mathrm{ref}})
\, L_{\mathrm{ref}},
\label{eq:criterion}
\end{equation}
where $C_{\mathrm{thr}}$ is the laboratory contrast threshold at the
prevailing adaptation, represented by Crumey's 2014 analytic
formulation of the Blackwell 1946 threshold-versus-intensity data
\citep{blackwell1946,crumey2014} and $f$ is the
field factor of Section~\ref{sec:calibration}. Fajr al-sadiq is declared
at the depression where the test passes simultaneously across the
band's lateral extent (the spread condition); passage of the central
patches alone, without the spread, is classified as the false dawn
(Section~\ref{sec:kadhib}). We refer to this test,
Eq.~\eqref{eq:criterion} together with the spread condition, as the
khayt (thread) criterion. The evening events mirror the test as
disappearances, with the red channel additionally gated at the cone
threshold: below approximately $10^{-3}$~cd\,m$^{-2}$, rod vision carries
no color, so there is no red shafaq left to lose; this gate matters
for the Hanafi/Shafi'i boundary at high latitude.

Because Eq.~\eqref{eq:criterion} is a ratio test against the same
night's sky, multiplicative systematics (absolute radiometric scale,
uniform cloud attenuation, uniform skyglow lift) largely cancel, and the
criterion remains defined wherever the sky physically brightens. At
$54.8^\circ$N midsummer, inside persistent twilight where no fixed-angle
convention returns a time at all, the model still returns a defined khayt
dawn (02:55 on the example date), illustrating that the relative
criterion stays defined where absolute conventions fail; this is a
demonstration of definedness, not a ground-truth validation point.

Figure~\ref{fig:night} shows the criterion applied to one example
night: the three detection margins per scanned depression, with the
false dawn reported where the central margin crosses threshold before
the spread margin.

\begin{figure}[t]
\centering
\includegraphics[width=0.72\textwidth]{p2_example_night}
\caption{One example night (Mecca, clear sky) is shown: the three
detection margins of Eq.~\eqref{eq:criterion} against the scanned solar
zenith angle. The horizontal line marks the detection threshold; the
gap between the central crossing and the spread crossing is the
false-dawn interval.}
\label{fig:night}
\end{figure}

Clouds require one further textual decision. The ayah's ``to you''
makes the criterion observer-relative: a dawn hidden by an eastern cloud
bank has not become distinct to the observer beneath it. The model
therefore evaluates the criterion through the measured cloud field
rather than above it; a regression test pins the expected asymmetries
(an eastern bank delays the dawn detection; a western bank advances the
evening disappearance, since dimming a fading band extinguishes its
distinctness earlier).

\section{Calibration and the universality of the field factor}
\label{sec:calibration}

The single calibrated constant $f$ in Eq.~\eqref{eq:criterion} absorbs
the laboratory-to-field gap of Section~\ref{sec:sources}. Fixing it
requires care, because the published desert campaigns do not report a
uniform statistic. The observing groups state their convention
explicitly: a bare headline depression is the ``highest value of
confidence,'' a mean or median plus one or two standard deviations, not
the central tendency \citep{hassan2014,hassan2015tubruq,hassan2016,semeida2018}.
Calibrating to those upper bounds as though they were central means --
as an earlier revision did, across the desert body as a whole -- biases
the target. We therefore fix $f$ only on the campaigns that report a
documented central mean, hold it worldwide, and read every other
campaign as a test. A decomposition of the calibrated value against the
published psychophysics literature (field factor for non-optimal
stimuli, blurred-edge degradation, binocular and practiced-observer
corrections, visibility-level practice in applied lighting) brackets it
comfortably; the decomposition is recorded in the validation archive of
the code release.

\begin{figure}[t]
\centering
\includegraphics[width=0.7\textwidth]{p2_edge_factor}
\caption{The per-site inversion of the field factor, from exact
fine-ladder runs (each point is the $f$ that a campaign implies through
the engine). The constant is pinned only by the campaigns with a
documented central mean (filled circles): Aswan, Riyadh, and Kottamia
imply 51, 53, and 57, an optimum of 53 with a clean-mean plateau to 60
(shaded); the production constant (green, 56) sits within it. Hail (red
diamond, implied $f=81$) is a central-mean outlier, its dawn
$14.01^{\circ}$ a half-degree shallower than any other clean mean; its
leverage alone pulled an earlier three-site fit to 70. The remaining
campaigns (open squares) report upper bounds or cross-study
photoelectric values, not central means; their implied factors happen
to fall in the same band because a site's mean\,$+$\,2SD dawn is its
clear-sky detection limit, which is what the engine predicts. Inverting
those campaigns' \emph{raw} central means instead scatters $f$ from 57
to 249 (off scale) and pins nothing.}
\label{fig:edge}
\end{figure}

The constant is pinned by inversion (Figure~\ref{fig:edge}): running a
campaign through the engine at a fine factor ladder and solving for the
$f$ that reproduces its dawn. Restricted to the campaigns with a
documented central mean -- Riyadh ($14.58^\circ$, naked eye), Aswan
($14.90^\circ$, calibrated camera), and Kottamia ($14.665^\circ$, naked
eye) -- the implied factors cluster tightly at 53, 51, and 57, with an
RMS-optimal $f = 53$ (residual $0.05^\circ$), the three agreeing to
within four units of $f$. Hail's central mean, $14.01^\circ$, is a
half-degree shallower than any other clean mean and implies
$f\approx81$; it is a genuine shallow-dawn outlier, and in an earlier
three-site Mecca cluster its leverage alone pulled the fit to 70.
Admitting Hail moves the clean optimum only to 60. The production
constant is set to $f = 56$, centrally within the $53$--$60$ band that
the central-mean campaigns allow, and stable to sub-degree shifts of
that basis.

The remaining desert campaigns do not report central means, and must
not be read as if they did. Their headline depressions are upper bounds
(Tubruq $14.7^\circ =$ mean\,$+$\,2SD, Sinai $14.61^\circ =$
median\,$+\,\sigma$, Assiut $13.665^\circ =$ median\,$+\,2\sigma$, Wadi
Al-Natrun $14.57^\circ =$ mean\,$+$\,1SD), while the Bahariya and
Matrouh figures come from a separate photoelectric study rather than the
naked-eye mean. Their raw central means are far shallower (Tubruq
$13.14^\circ$, Assiut $11.25^\circ$, Matrouh $13.41^\circ$), contaminated
by all-nights averaging and light pollution that the clear-sky engine
does not model; inverting those raw means scatters the implied factor
from 57 to 249 and pins nothing. That the engine at $f = 56$ nonetheless
reproduces the \emph{upper bounds} is physical, not fitted: a site's
mean\,$+$\,2SD dawn is its deepest, best-visibility, clear-sky detection
-- exactly the regime the engine predicts. The clean central means and
the contaminated campaigns' upper bounds therefore coincide near
$14.5$--$14.9^\circ$ because both are the clear-sky detection limit,
seen two ways. The inversion shows no significant trend with latitude.

The calibration was executed on the frozen final engine by exact GPU
runs at a fine factor ladder (the tool and per-site runs are archived
with the code release). All campaign numbers in this paper are quoted at
$f = 56$; every site outside the three central-mean anchors is a genuine
test. The Birmingham out-of-sample statistics are insensitive to the
choice within the $53$--$60$ plateau. A later deep-cloud estimator
upgrade and an evening-solver bracket correction (companion paper and
code history) leave the criterion's dawn detection depressions unchanged
to within $0.02^{\circ}$ at the calibration anchors, so the calibration
is insensitive to both.

\section{Validation against observation campaigns}
\label{sec:campaigns}

\begin{table}[t]
\centering
\caption{Model versus published dawn campaigns (abbreviated; the full
table with per-site run brackets is in the repository validation
documents). ``Engine'' is the khayt criterion output matched to the
campaign's modality (naked eye versus instrument). ``Basis'' is the
statistic the source reports for its headline depression, which is
\emph{not} uniform across campaigns: several are documented upper
bounds (mean or median $+\,k\sigma$, the group's stated ``highest value
of confidence''), not central means. Residuals are engine minus the
published headline value, in degrees of solar depression; for the
upper-bound rows this is a lower bound on the mean-basis agreement
(the engine sits at the clear-sky detection limit, which coincides
with those sites' best-visibility tails). The constant $f$ is pinned
on the central-mean campaigns only (Section~\ref{sec:calibration});
every row is an out-of-sample test at fixed $f$.}
\label{tab:campaigns}
\footnotesize
\setlength{\tabcolsep}{3pt}
\begin{tabular}{lcclc}
\toprule
Site (campaign) & Engine & Observed & Basis & Residual \\
\midrule
Riyadh desert (KACST) & 14.50 & $14.6 \pm 0.3$ & mean & $-0.10$ \\
Hail & 14.46 & $14.01 \pm 0.32$ & mean & $+0.45$ \\
Aswan 2016, camera & 14.82 & $14.90 \pm 0.17$ & camera mean & $-0.08$ \\
North Sinai & 14.16 & 14.61 & med$+\sigma$ & $-0.45$ \\
Fayum & 14.41 & 14.8 / SQM 13.75--14.7 & range & in range \\
Matrouh & 14.32 & 14.5 & photoelec.\ & $-0.18$ \\
Kottamia & 14.75 & $14.66 \pm 0.2$ & mean & $+0.09$ \\
Bahariya & 14.63 & 14.6 & photoelec.\ & $+0.03$ \\
Wadi Al-Natrun & 14.00 & range 12.48--15.14 & mean$+$1SD & inside \\
Tubruq (desert horizon) & 14.68 & 14.66--14.7 & mean$+$2SD & $0.00$ \\
Tubruq (sea horizon) & 14.68 & 13.43--13.48 & mean$+$2SD & $+1.22$ \\
Assiut (agricultural) & 14.69 & 13.665 & med$+2\sigma$ & $+1.03$ \\
Sana'a, 2200 m & 15.04 & bracket [13.2, 18.8] & bracket & inside \\
Depok, SQM knee & 14.23 & $14.0 \pm 0.6$ & SQM knee & $+0.23$ \\
Malaysia SQM sites & 14.48 & $14.19 \pm 0.52$ & SQM knee & $+0.29$ \\
Aswan 1909, camera & 14.82 & 14.25 & camera & $+0.57$ \\
\bottomrule
\multicolumn{5}{p{0.95\columnwidth}}{\footnotesize ``mean/med $+\,k\sigma$'' rows are
published upper bounds; the sites' central means are shallower (e.g.\
Tubruq desert $13.14^\circ$, Assiut $11.25^\circ$) but reflect
all-nights or light-polluted averages the clear-sky engine does not
model. ``photoelec.'' denotes a cross-study photoelectric value, not
the naked-eye mean (Bahariya $13.81^\circ$, Matrouh $13.41^\circ$).
$f$ is pinned on the central-mean campaigns (Riyadh, Aswan, Kottamia),
not fitted per row.}
\end{tabular}
\end{table}

\begin{figure}[t]
\centering
\includegraphics[width=0.8\textwidth]{p2_site_map}
\caption{The validation campaign sites are mapped, with the inset
expanding the Egyptian and Libyan cluster. Mecca (starred) anchors the
criterion historically; the constant is pinned on the central-mean
campaigns only (Riyadh, Aswan, Kottamia), so every other campaign is an
out-of-sample test.}
\label{fig:map}
\end{figure}

\begin{figure}[t]
\centering
\includegraphics[width=0.78\textwidth]{p2_campaign_forest}
\caption{Residuals of Table~\ref{tab:campaigns} are shown as a forest
plot (engine minus observed, with the campaigns' published uncertainties;
the shaded band marks $\pm 0.3^{\circ}$, close to the pooled
median absolute residual of $0.26^{\circ}$). The two documented open misses, both background-sky
physics rather than criterion failures, are discussed in the text.}
\label{fig:forest}
\end{figure}

Figure~\ref{fig:map} maps the campaign sites,
Table~\ref{tab:campaigns} summarizes the campaign comparison, and
Figure~\ref{fig:forest} shows the residuals with their published
uncertainties. Because eye and instrument measure different events
(Section~\ref{sec:sources}), the residual statistics are reported per
modality: over the naked-eye campaigns the median absolute residual is
$0.28^{\circ}$, and over the instrumental campaigns (SQM knees and
calibrated cameras) $0.26^{\circ}$; pooled over all fourteen scored rows
it is $0.26^{\circ}$. The constant $f$ is pinned on the three
central-mean campaigns (Section~\ref{sec:calibration}), and shifting
that basis moves it by less than a degree, so each remaining campaign
functions as an out-of-sample test rather than a fitted point. For the
rows whose published depression is an upper bound
(Table~\ref{tab:campaigns}), the tabulated residual is a lower bound on
the mean-basis agreement; it stays small because the engine sits at the
clear-sky detection limit those bounds also mark. The Sana'a row is
scored only as a consistency check: Sultan's two bracketing events
(first glow merged with the zodiacal light at $18.8^{\circ}$; colour
divergence at $13.2^{\circ}$) straddle any plausible detection value,
so the row cannot discriminate between models and is excluded from the medians. This
median of about a quarter-degree of depression is
roughly a minute and a half of clock time at
mid-latitudes, and comparable to the scatter between the campaigns
themselves. Four features deserve comment.

First, the instrument-versus-eye split behaves as the psychophysics
requires: the model emits a distinct instrumental (absolute-threshold)
output and a naked-eye (khayt) output that differ by about half a degree
of depression, and each is compared against data of its own modality. The
instrumental output matches the SQM campaigns at the residuals tabulated
in Table~\ref{tab:campaigns} ($+0.23^{\circ}$ at Depok, $+0.29^{\circ}$
for the Malaysia sites), while the khayt output matches the visual
panels. A 116-year-old photographic campaign \citep{miethe1909} is
reproduced at $+0.57$.

Second, one apparent outlier was resolved at its primary source: the
Wadi Al-Natrun target circulating in secondary summaries ($14.57^{\circ}$)
is that paper's mean-plus-one-standard-deviation upper confidence
bound, not its campaign mean; the published observed range is
$12.48$ to $15.14^{\circ}$ over 38 mornings, and the model sits inside
it. We flag this as a general hazard of campaign meta-analysis: the
Egyptian and Saudi campaign literature quotes upper-confidence bounds
(``white thread'' safety values) alongside means, and a comparison must
match the statistic, not the headline number.

Third, the remaining misses are physical rather than statistical, and
they are documented as such.
The two large residuals sit exactly where the reference-sky assumption
breaks: Assiut ($+1.03$), where the observers faced an agricultural
valley whose aerosol and humidity regime the desert climatology cannot
represent (the aerosol channel of Section~\ref{sec:inputs} is the
identified mechanism, feeding measured optical depth moving the residual
back toward the observed range); and the Tubruq sea horizon
($+1.22$), where the dawn path crosses a directional marine boundary
layer that no one-dimensional aerosol column can represent. For the
sea-horizon row that mechanism has now been tested directly: representing
the marine boundary layer as a georeferenced haze slab confined to the
sea cells along the dawn path, at climatological marine optical depth
(0.12 at 550\,nm, 700\,m scale height), shifts the computed sea-horizon
dawn by $-0.54^{\circ}$ of depression averaged over three seasons,
closing roughly half
the residual (from $+1.2$ to about $+0.5^{\circ}$), while the same
observers' desert horizon remains matched at $0.00$. The remaining half
is kept as a documented open row rather than absorbed into the
calibration; the candidate explanations are marine haze episodically
thicker than the climatology, and a different visual task over a
featureless sea horizon.

\begin{figure}[t]
\centering
\includegraphics[width=0.78\textwidth]{p2_birmingham}
\caption{The Birmingham year. OpenFajr consensus-panel dawn
depressions (black; all 42 published dates) are compared with the
engine under clear sky with zero retuning (blue) and under the full
measured atlas veil (purple). Both peaks of the double-peaked curve
and the solstice troughs are present in the engine's own seasonal
structure (statistics in the text). The winter dates carry the urban
veil discussed in Section~\ref{sec:inputs}; with the measured skyglow
they close to within the reported uncertainty band.}
\label{fig:bham}
\end{figure}

Fourth, the decisive out-of-sample test is Birmingham
(Figure~\ref{fig:bham}). The OpenFajr
program \citep{openfajr2016} imaged 300 mornings over a full year at
$52.4^\circ$N and had a 19-member panel (professional astronomers,
timekeeping authorities, and experienced mosque observers) judge 42
clear-horizon date series, producing a double-peaked seasonal curve of
dawn depressions between 12.3 and 15.0. The model, calibrated on desert
campaigns some twenty degrees of latitude away and never shown these
data, reproduces the summer
trough (computed 11.9 to $12.6^\circ$ across the eight June dates
after the compressed-twilight scan refinement of Paper~I; panel 12.3
to 12.7) and both peaks of the seasonal curve. The summer relaxation
that UK scholars apply as an empirical rule (from $\sim 14.5^\circ$ in
winter toward $\sim 12.5^\circ$ in midsummer) emerges from
persistent-twilight geometry with no additional freedom. Across all
42 dates the clear-sky statistics are mean $+0.54^\circ$, RMS
$0.95^\circ$, $r = 0.51$, and the winter share of the residual is
urban-veil physics rather than criterion error: running every winter
date under the full measured atlas veil brackets the panel from the
other side (mean $-1.13^\circ$, and the panel lies between the
clear-sky and full-veil runs on 13 of the 14 winter dates). The
bracket width is the quantitative statement of the one input no
public feed supplies, the street-light duty cycle at the dawn hour;
the panel's position inside it corresponds to a dawn-hour veil of
roughly half the atlas's overpass-time value, consistent with
documented UK part-night dimming practice
(Section~\ref{sec:inputs}).

\subsection{The evening events}
\label{sec:isha}

The evening disappearances rest on thinner data than dawn, and the
comparison is correspondingly shorter (Table~\ref{tab:isha}). The
white-twilight end lies within $0.66^{\circ}$ of the instrumental
Mecca-region value ($17.33$ versus $17.99 \pm 0.16^{\circ}$) and
within $0.33^{\circ}$ of the classical muwaqqit convention of
$17^{\circ}$. The engine's red-end value ($14.96^{\circ}$) sits at the
upper bound of the sparse visual literature (12
to $15^{\circ}$), essentially coincident with Sultan's $15^{\circ}$. No
color-resolved modern
campaign exists; until one does, the red calibration should be read
as provisional.

\begin{table}[t]
\centering
\caption{Evening (isha) validation. The white event is matched to
instrumental twilight-end data; the red event to the visual
literature and Sultan's observation.}
\label{tab:isha}
\footnotesize
\setlength{\tabcolsep}{4pt}
\begin{tabular}{lcc}
\toprule
Event & Engine (Mecca) & Independent value \\
\midrule
Shafaq al-abyad (white end) & $17.33^{\circ}$ & SQM end $17.99 \pm 0.16^{\circ}$; muwaqqit $17^{\circ}$ \\
Shafaq al-ahmar (red end) & $14.96^{\circ}$ & literature 12 to $15^{\circ}$; Sultan $15^{\circ}$ \\
\bottomrule
\end{tabular}
\end{table}

\subsection{Against the conventions in use}

The same 42 panel dates measure the methods the world actually uses
(Table~\ref{tab:conventions}). The Muslim World League angle of
$18^{\circ}$ does not exist on 20 of the 42 dates (the Sun never
reaches that depression at $52.4^{\circ}$N in summer) and misses by
more than four degrees, roughly half an hour of clock time, where it
does exist. The $15^{\circ}$ convention exists everywhere but misses
by $1.76^{\circ}$ RMS. The physics-based model is defined on every date of
the year at every latitude. Over all 42 dates it is a factor of two more
accurate than the best year-round fixed angle (ISNA $15^{\circ}$,
$1.76^{\circ}$ against the model's $0.95^{\circ}$ RMS); the most widely
used angle ($18^{\circ}$) is undefined on 20 of the 42 dates and misses
by $4.4^{\circ}$ RMS on the 22 where it exists.

\begin{table}[t]
\centering
\caption{The OpenFajr panel year as a benchmark of methods: residuals
of each method against the 42 consensus-panel dates.}
\label{tab:conventions}
\small
\begin{tabular}{lccc}
\toprule
Method & Mean [deg] & RMS [deg] & Dates defined \\
\midrule
This work, clear-sky model & $+0.54$ & $0.95$ & 42 of 42 \\
Fixed $15^{\circ}$ (ISNA) & $+1.62$ & $1.76$ & 42 of 42 \\
Fixed $18^{\circ}$ (MWL) & $+4.32$ & $4.38$ & 22 of 42 \\
\bottomrule
\end{tabular}
\end{table}

\section{The false dawn}
\label{sec:kadhib}

Because the criterion evaluates the band patches individually, the false
dawn requires no separate machinery: when the central patches pass the
distinctness test while the spread condition still fails, the model
reports fajr al-kadhib, the narrow vertical brightening of the zodiacal
column. At the calibration sites the computed kadhib precedes the sadiq
by roughly $0.3$ to $0.7^\circ$ of depression (one and a half to three
and a half minutes, varying with site and season), consistent with the
classical description of a brief interval in which eating remains
permitted while the cautious wait.

\begin{figure}[t]
\centering
\includegraphics[width=0.6\textwidth]{p2_kadhib}
\caption{Computed suppression of the false dawn's distinctness by
artificial skyglow. The lead of the kadhib detection over the sadiq
is shown, in degrees of depression, as the veil is swept from a
clear sky to Bortle class 6 at fixed geometry. The monotone
compression reproduces the qualitative visibility-literature
expectation on a phenomenon that the calibration never used.}
\label{fig:kadhib}
\end{figure}

The zodiacal character of the event provides an independent
test on which the constant was never tuned: the visibility
literature places the loss of the zodiacal light near Bortle class 4 to
5 skies. When the artificial veil is swept at fixed geometry, the
computed kadhib lead over the sadiq compresses monotonically from
$0.48^\circ$ under a clear sky to $0.33^\circ$ at Bortle 6
(Figure~\ref{fig:kadhib}), reproducing the qualitative suppression; the model does not reproduce a
hard extinction of the event, which we state as a partial confirmation,
since the computed kadhib occurs immediately ahead of dawn, where the
wedge is brightest, a configuration the free-standing-cone literature
does not directly address.

\section{Environmental response and reported uncertainty}
\label{sec:inputs}

Because the criterion is evaluated on tonight's simulated sky, the model
quantifies how the times respond to measured conditions, and it reports
a propagated uncertainty with every time.

\textbf{Clouds.} A measured three-dimensional cloud field
(satellite-reconstructed; Paper~I) shifts the dawn detection with the
azimuthal asymmetry the observer-relative reading requires: an eastern
stratus bank delays fajr while leaving isha nearly untouched, and a
western bank advances the evening disappearance. Uniform thin cirrus
principally dims the band and delays detection by minutes; the criterion's
ratio structure cancels much of a uniform deck's effect, which is why
overcast mornings remain computable at all.

\textbf{Aerosol.} The model ingests the Copernicus Atmosphere
Monitoring Service (CAMS) aerosol optical depth at
the prayer hour as an excess over a baseline measured at the desert
calibration campaigns (chosen so the clear-desert body itself carries
zero excess by construction). The measured response is 13 to 17 degrees of depression
per unit excess optical depth in the linear regime, with the saturation
behavior the crossing slope predicts. A sensitivity run at the shipped
constant confirms the direction and scale: a continental-average aerosol
load representative of the agricultural regime reduces the predicted
Assiut depression by of order $1.5^\circ$, enough to carry the residual
across the observed value, while an urban load overcorrects, so the
measured optical depth is bracketed between the clean desert baseline and
a continental one. The desert calibration sites carry near-zero measured
excess by construction and are therefore essentially unmoved, the
differential response expected when one site sits in an aerosol regime
the desert baseline does not carry.

\textbf{Moonlight.} The Moon enters the criterion on both sides:
moonlight raises the band patches and the reference patches together,
and because the test is a ratio against the same night's sky, a bright
Moon principally raises the adaptation level (through the
threshold-versus-intensity function) rather than shifting the
detection wholesale. The lunar state is taken from the ephemeris
(true phase angle, parallax, and distance), and the OpenFajr data
provide an implicit end-to-end check: the panel reported no lunar
structure in their year of detections, and none of the model's
residuals against the 42 dates correlates with lunar phase.

\textbf{Artificial light.} The veil is taken from the propagated world
atlas, cross-checked against Visible Infrared Imaging Radiometer Suite
(VIIRS) nighttime lights, and applied in mesopic photometry
computed from the source spectrum: a sodium-dominated city veils
threshold vision at barely a third of its photopic luminance, an LED
city at nearly all of it, a distinction worth more than a factor of two
at the adaptation levels of dawn observation. Where part-night switching
makes the veil bimodal, the model widens the reported uncertainty
accordingly.

\begin{figure}[t]
\centering
\includegraphics[width=\textwidth]{p2_aod}
\caption{The left panel shows the khayt dawn depression versus the
measured excess aerosol optical depth for the validation sites (each
point is one campaign date run with the CAMS value at the prayer
hour). The right panel shows the propagated uncertainty of each run
split by source; the Birmingham winter row is dominated by the skyglow
term, itself dominated by the street-light duty cycle.}
\label{fig:aod}
\end{figure}

Figure~\ref{fig:aod} shows the measured-aerosol response and the
per-run uncertainty decomposition.

\textbf{Reported uncertainty.} Each computed time carries a standard
deviation combining Monte Carlo noise, the atlas calibration, the
duty-cycle term, and the aerosol input, each propagated through the
measured sensitivity of the detection margin at the crossing.
Representative values: Mecca $\pm 3.0$ minutes; a clear coastal site
$\pm 1.8$; Birmingham in winter $\pm 7.2$, dominated by the duty-cycle
term. We regard the width of the urban winter band as a finding
in itself: it is a quantitative statement of how much a calendar can
currently know.

\section{Discussion}
\label{sec:discussion}

\textbf{Jurisprudential scope.} The model computes the physical
observable that the sources describe; it does not adjudicate between
legal schools. Where the schools differ, the model serves each on its own terms
rather than dictating between them: for the isha boundary it tracks
the red and the white disappearances as separate physical events
(shafaq al-ahmar for the majority Maliki, Shafi'i, and Hanbali
position; shafaq al-abyad for the Hanafi school), reporting both
times on every night, and the physically forced gap between them
(the red band vanishes at the cone threshold while the white persists)
gives each school an objective observable matched to its own
definition, and its high-latitude behavior
(times defined by relative brightening) is offered as physical input to
a juristic discussion, not as a ruling. Computed times are explicitly
experimental for worship and should be cross-checked against established
local calendars; the appropriate near-term use is as an instrument: for
calendar committees evaluating their conventions against physics, for
campaign design, and for quantifying the environmental corrections
(cloud, aerosol, skyglow) that no fixed angle can carry.

\textbf{Operational accessibility.} Because the model's accuracy
rests on measured inputs, its practical use by a mosque committee or
observatory does not require them: every dynamic input degrades to a
documented static fallback (OPAC climatological aerosol by region,
the atlas value or a Bortle class for skyglow, clear sky when no
cloud product is available, the built-in solar ephemeris), and the
computed time then simply carries a wider reported uncertainty, so a
committee can run the model offline from coordinates and a date alone
and see, from the widths, exactly which measured input would tighten
their calendar most. Precomputing an annual table at climatological
inputs, as conventional timetables do, is a single command.

\textbf{Relation to prior computational work.} Prior quantitative
treatments of fajr visibility have combined twilight brightness models
of the Rozenberg lineage with threshold curves \citep{ilyas1984}, or fit
depression statistics directly. The present contribution is, to our
knowledge, the first to compute the criterion from absolute spectral
radiative transfer validated against reference codes and measured skies
(Paper~I), the first to carry measured three-dimensional cloud fields,
measured aerosol, and spectrally resolved artificial light into the
detection, and the first to report per-event propagated uncertainties.

\textbf{Boundaries of the present validation.} The evening red event rests on the weakest data
(no color-resolved modern campaign exists; the calibration leans on
instrumental values and on Sultan's bracketing observations). The
sea-horizon problem is open. The field factor, though cross-site invariant,
is a single scalar summarizing a rich psychophysical literature; a
purpose-built field campaign with controlled adaptation and modern
threshold methodology would either confirm it or replace it with
something better, and the release includes the instrument protocol for
exactly that campaign. Finally, the transport itself carries the
documented deep-cloud statistical limits of Paper~I.

\section{Conclusions}
\label{sec:conclusion}

The Quranic dawn criterion is a well-posed detection problem, and
solving it as one reproduces more than a century of dawn observation
with a median absolute residual of $0.26^\circ$ of depression (pooled over
all scored campaigns) from a
single calibrated constant whose per-site inversion shows a cross-site
spread the size of the campaigns' own noise with no latitude trend. A
fixed angle is the correct first-order summary of dark-desert dawn
(the campaigns cluster near $14$ to $15^\circ$), and the model explains
both why that summary works where it works and how it fails where it
fails: with season and latitude (Birmingham's curve), with aerosol
regime (Assiut), with background sky (Tubruq's two horizons), and with
urban light (the winter veil). What replaces the angle is not another
constant but a computation, carrying measured inputs and explicitly
propagated uncertainties, offered to observers, calendar committees, and jurists as
an instrument for the discussion that remains theirs.

\appendix

\section{Birmingham per-date comparison}

\begin{table}[h]
\centering
\caption{All 42 OpenFajr panel dates at the recalibrated constant.
Panel values are the 19-member consensus depressions; the clear-sky
residual is engine (clear sky) minus panel; the final column is the
engine under the full measured atlas veil, which together with the
clear-sky column brackets the panel (Section~\ref{sec:campaigns}).}
\label{tab:bham}
\scriptsize
\setlength{\tabcolsep}{4pt}
\begin{tabular}{lcccc}
\toprule
Date 2015 & Panel & Clear sky & Residual & Full veil \\
\midrule
Jan 11 & 13.0 & 13.83 & $+0.83$ & 12.05 \\
Jan 19 & 13.1 & 14.83 & $+1.73$ & 12.07 \\
Jan 24 & 12.9 & 14.59 & $+1.69$ & 12.03 \\
Feb 06 & 13.7 & 13.68 & $-0.02$ & 12.04 \\
Feb 18 & 13.6 & 14.61 & $+1.01$ & 12.06 \\
Feb 22 & 13.7 & 14.65 & $+0.95$ & 12.06 \\
Feb 23 & 13.2 & 14.61 & $+1.41$ & 12.07 \\
Feb 24 & 12.9 & 14.69 & $+1.79$ & 12.06 \\
Feb 27 & 13.0 & 14.48 & $+1.48$ & 12.07 \\
Apr 20 & 15.0 & 14.53 & $-0.47$ & 12.07 \\
Apr 22 & 14.5 & 14.70 & $+0.20$ & 12.07 \\
Apr 27 & 13.7 & 14.47 & $+0.77$ & 12.07 \\
Apr 28 & 13.6 & 14.56 & $+0.96$ & 12.06 \\
May 13 & 13.0 & 14.59 & $+1.59$ & 12.07 \\
May 17 & 12.8 & 14.67 & $+1.87$ & 12.06 \\
May 20 & 13.1 & 14.51 & $+1.41$ & 12.06 \\
May 23 & 13.0 & 14.16 & $+1.16$ & 12.04 \\
May 27 & 13.0 & 13.74 & $+0.74$ & 12.04 \\
Jun 04 & 12.9 & 12.56 & $-0.34$ & 11.74 \\
Jun 06 & 12.7 & 12.58 & $-0.12$ & 11.72 \\
Jun 07 & 12.6 & 12.58 & $-0.02$ & 11.74 \\
Jun 08 & 12.7 & 12.58 & $-0.12$ & 11.74 \\
Jun 09 & 12.7 & 12.61 & $-0.09$ & 11.74 \\
Jun 10 & 12.6 & 12.64 & $+0.04$ & 11.75 \\
Jun 22 & 12.5 & 11.90 & $-0.60$ & 11.36 \\
Jun 30 & 12.3 & 12.60 & $+0.30$ & 11.73 \\
Jul 06 & 13.0 & 12.57 & $-0.43$ & 11.74 \\
Jul 10 & 13.0 & 12.69 & $-0.31$ & 11.76 \\
Jul 18 & 13.8 & 14.12 & $+0.32$ & 12.04 \\
Jul 21 & 13.7 & 13.98 & $+0.28$ & 12.04 \\
Aug 16 & 14.3 & 14.79 & $+0.49$ & 12.07 \\
Aug 28 & 14.2 & 13.99 & $-0.21$ & 12.03 \\
Sep 06 & 14.9 & 14.18 & $-0.72$ & 12.05 \\
Sep 07 & 14.5 & 14.35 & $-0.15$ & 12.06 \\
Sep 23 & 14.6 & 14.65 & $+0.05$ & 12.06 \\
Sep 26 & 14.4 & 14.56 & $+0.16$ & 12.06 \\
Sep 27 & 14.3 & 13.67 & $-0.63$ & 12.03 \\
Nov 13 & 13.5 & 14.76 & $+1.26$ & 12.02 \\
Nov 28 & 13.4 & 13.67 & $+0.27$ & 12.03 \\
Dec 10 & 12.9 & 14.60 & $+1.70$ & 12.07 \\
Dec 19 & 13.1 & 14.62 & $+1.52$ & 12.06 \\
Dec 25 & 12.6 & 13.57 & $+0.97$ & 12.03 \\
\bottomrule
\end{tabular}
\end{table}

The winter rows of Table~\ref{tab:bham} carry the site's urban veil.
At the recalibrated constant the clear-sky 42-date statistics are mean
$+0.54^{\circ}$, RMS $0.95^{\circ}$; the full-veil column brackets
the panel from the other side, and the azimuthally resolved veil
(slant-path Garstang on the measured VIIRS radiance field,
Section~\ref{sec:inputs}) narrows the bracket from the veiled side
(January 24: 11.82 versus the panel's 12.9; December 10: 11.79 versus
12.9). The residual separating the directional-veil runs from the
panel is the street-light duty cycle at the dawn hour, the one input
no public feed supplies; the June dates are compared clear-sky because
their 01:22 GMT dawns fall inside the documented lights-dimmed
window.

\section*{Code and data availability}

The simulator and the detection layer are available at
\url{https://github.com/Muno459/twilight} (MIT/Apache-2.0); release
v1.0.0 corresponds to this paper. The release includes a validation
archive containing the campaign run records and the exact data behind
every figure and table, together with the scripts that regenerate
them. The companion methods paper (Paper~I) documents the radiative
transfer validation.

\section*{Funding and competing interests}

This research received no external funding. The author declares no
competing interests.

\section*{Acknowledgements}

This work stands on the campaigns cited in Table~\ref{tab:campaigns}:
observers who rose before dawn, month after month, from Aswan in 1908 to
Birmingham in 2015, and published what they saw.

\bibliographystyle{elsarticle-harv}
\begin{thebibliography}{00}

\bibitem[Al-Mostafa et al.(2005)]{almostafa2005}
Al-Mostafa, Z.A., et al., 2005. Studying of Twilight Project, first
part. Technical report, Institute of Astronomical and Geophysical
Research, King Abdulaziz City for Science and Technology (KACST),
Riyadh. One-year naked-eye and CCD morning-twilight campaign, deep
desert site 170~km from Riyadh.

\bibitem[Blackwell(1946)]{blackwell1946}
Blackwell, H.R., 1946. Contrast thresholds of the human eye. Journal of
the Optical Society of America 36, 624--643.

\bibitem[Crumey(2014)]{crumey2014}
Crumey, A., 2014. Human contrast threshold and astronomical visibility.
Monthly Notices of the Royal Astronomical Society 442, 2600--2619.

\bibitem[Hassan et al.(2014)]{hassan2014}
Hassan, A.H., Abdel-Hadi, Y.A., Issa, I.A., Hassanin, N.Y., 2014.
Naked-eye observations for morning twilight at different sites in Egypt.
NRIAG Journal of Astronomy and Geophysics 3, 23--26. Reports a pooled
grand mean $13.642^\circ \pm 1.054^\circ$ ($n=35$) and a headline
``$14.7^\circ$'' quoted explicitly as mean\,$+$\,1SD.

\bibitem[Hassan and Abdel-Hadi(2015)]{hassan2015tubruq}
Hassan, A.H., Abdel-Hadi, Y.A., 2015. Determination of the beginning of
morning twilight at Tubruq, Libya. Middle-East Journal of Scientific
Research 23, 2627--2632. Desert-site mean $13.144^\circ \pm 0.757^\circ$
($n=623$); the quoted $14.7^\circ$ is mean\,$+$\,2SD.

\bibitem[Hassan et al.(2016)]{hassan2016}
Hassan, A.H., Issa, I.A., Mousa, A.M., Abdel-Hadi, Y.A., 2016. Naked-eye
determination of the dawn for Sinai and Assiut of Egypt. NRIAG Journal
of Astronomy and Geophysics 5, 9--15. Sinai $14.61^\circ$
(median\,$+\,\sigma$); Assiut $13.665^\circ$ (median\,$+\,2\sigma$;
central mean $11.25^\circ$).

\bibitem[Ilyas(1984)]{ilyas1984}
Ilyas, M., 1984. A Modern Guide to Astronomical Calculations of Islamic
Calendar, Times and Qibla. Berita Publishing, Kuala Lumpur.

\bibitem[Semeida and Hassan(2018)]{semeida2018}
Semeida, M.A., Hassan, A.H., 2018. Determination of the beginning of
morning twilight using the naked eye at Wadi Al-Natrun, Egypt.
Beni-Suef University Journal of Basic and Applied Sciences 7, 286--290.
Headline $14.57^\circ$ stated as mean\,$+$\,1SD (``highest value of
confidence''); observed range $12.48$--$15.14^\circ$ ($n=38$).

\bibitem[Khalifa et al.(2018)]{khalifa2018}
Khalifa, N.S., Hassan, I.A., Taha, M.R., 2018. Naked-eye determination
of the solar depression angle of true dawn at Hail, Saudi Arabia. NRIAG
Journal of Astronomy and Geophysics 7, 22--26.

\bibitem[Mawad(2024)]{mawad2024}
Mawad, R., 2024. The height of the diurnal atmosphere: twilight
altitude. Advances in Space Research 74.

\bibitem[Miethe and Lehmann(1909)]{miethe1909}
Miethe, A., Lehmann, E., 1909. D\"ammerungsbeobachtungen in Assuan,
Winter 1908. Meteorologische Zeitschrift 26.

\bibitem[OpenFajr(2016)]{openfajr2016}
OpenFajr Research Project, 2016. Research paper: an observational
study of the time of fajr in Birmingham, United Kingdom (May 2016,
with addendum September 2016).
\url{https://www.openfajr.org/docs/research_paper.pdf} (accessed
2026-07-06). All-sky CCD, 300 mornings; 42 clear-view series judged
by a 19-member consensus panel.

\bibitem[Sultan(2004)]{sultan2004}
Sultan, A.H., 2004. Fajr observations at Bani-Hoshiesh, Yemen.
al-Irshaad 8.

\end{thebibliography}

\end{document}
